\section{Radiomics and radiomics-related artificial intelligence on the road to clinical translation}

Majumder et al. (2024) presented a comprehensive state-of-the-art review discussing how radiomics and artificial intelligence (AI) are transforming oncologic imaging by extracting thousands of quantitative imaging features from conventional modalities such as CT, MRI, PET, and ultrasound \cite{Majumder2024-tzad004}. The authors described the full radiomics pipeline, beginning with standardized image acquisition and preprocessing steps such as resampling, normalization, discretization, and bias field correction, followed by segmentation using tools like ITK-SNAP or 3D Slicer. Radiomic features were grouped into four main categories: intensity-based, volumetric and morphologic, texture-based, and higher-order features derived through filtering or neural networks. To minimize site-specific variability, harmonization methods such as ComBat were emphasized as essential for multi-center reproducibility. Feature selection techniques including correlation analysis, LASSO regularization, and principal component analysis were used to reduce redundancy and improve generalizability. The review also highlighted the emergence of featureless deep learning (DL) methods, particularly convolutional neural networks (CNNs), which autonomously extract complex image representations and outperform traditional handcrafted approaches in segmentation, lesion detection, and prognostic prediction tasks. Various machine learning algorithms such as support vector machines, random forests, and logistic regression were reported for classification tasks, while Cox regression and Kaplan-Meier analyses were employed for survival modeling. The authors noted that validation typically involves cross-validation or external datasets and that performance metrics like AUC, sensitivity, and specificity are standard. Despite rapid progress, the study identified key barriers to clinical translation, including lack of data standardization, feature redundancy, segmentation variability, and interpretability of deep features. The authors concluded that future advances in automated segmentation, federated learning, and workflow harmonization will be critical for achieving reliable, explainable, and clinically deployable radiomic models in precision oncology.
